% ==========================================
% 2. Problem Definition
% ==========================================
%
% Status: stable.
% it fixes the common frame shared by all experiments.
% The research questions may be refined during the project, but the original
% hypotheses are never rewritten to fit the results (the gap is discussed in
% Section~\ref{sec:discussion}).
%
% Introduction vs this section: the introduction states the problem briefly;
% this section formalizes it. We avoid repeating the general context here.

\section{Problem Definition}
\label{sec:problem-definition}

% ------------------------------------------
% Short framing paragraph
% ------------------------------------------
% TODO [NOW]: We open with two or three sentences announcing the content of the
%             section (problem, questions, scope, evaluation criteria).

\subsection{Problem Statement}
\label{sec:problem-statement}
% TODO [NOW]: We formulate the problem precisely: what is given (inputs), what is
%             expected (outputs), and which properties the solution must satisfy.
% TODO [NOW]: We use a formal description (definitions, sets, functions) when it
%             removes ambiguity.

\subsection{Definitions and Notation}
\label{sec:notation}
% TODO [NOW]: We define once, here, every technical term and symbol used in the
%             report. All authors use these definitions, and no section
%             redefines them.
% TODO [DURING]: We add any new term as soon as an experiment introduces it.

\subsection{Challenges}
\label{sec:challenges}
% TODO [NOW]: We explain why the problem is hard, separating the technical
%             difficulties from the scientific ones (open questions).
% TODO [NOW]: We support each challenge with a reference or with a precise argument.

\subsection{Research Questions and Hypotheses}
\label{sec:rq-hypotheses}
% TODO [NOW]: We list the research questions, each one precise, answerable by an
%             experiment and numbered (RQ1, RQ2, ...). Each carries a label
%             \label{rq:1}, \label{rq:2}, ... so that other sections can refer to it.
% TODO [NOW]: For each research question, we state the hypothesis (H1, H2, ...)
%             and the reasoning behind it.
% TODO [DURING]: When a research question is refined, we keep a trace of the
%                change (date and reason) in notebook/.
% TODO [NOW]: We give an overview table linking each research question to its
%             hypothesis, the experiments addressing it and its metrics.
% Example structure:
%   \begin{table}[htbp]
%       \centering
%       \caption{Research questions, hypotheses and experiments.}
%       \label{tab:rq-overview}
%       \begin{tabular}{llll}
%           \toprule
%           RQ  & Hypothesis & Experiments & Metrics \\
%           \midrule
%           RQ1 & H1         & EXP-001     & ...     \\
%           \bottomrule
%       \end{tabular}
%   \end{table}
% TODO [CHECK]: Every research question is covered by at least one experiment,
%               and every experiment addresses at least one research question.

\subsection{Scope and Constraints}
\label{sec:scope}
% TODO [NOW]: We state what is in scope and what is explicitly out of scope,
%             with the reason for each exclusion.
% TODO [NOW]: We state the assumptions and constraints (time, hardware, data
%             availability, software) that frame the study.
% TODO [LAST]: We reread this subsection against the final experiments and
%              adjust it if the scope changed.

\subsection{Evaluation Criteria}
\label{sec:evaluation-criteria}
% TODO [NOW]: We define the metrics shared by all experiments (name, definition,
%             unit, and how it is computed), so that results are comparable.
% TODO [NOW]: We state what counts as a success or a failure for each research
%             question (thresholds, comparison baseline).
% TODO [NOW]: We state the statistical conventions: number of repetitions,
%             how variability is reported (standard deviation, confidence interval).
% TODO [DURING]: An experiment-specific metric is added here only if it is
%                defined in the same terms as the shared ones.
